\documentclass[12pt]{article}
\usepackage[utf8]{inputenc}
\usepackage[T1]{fontenc}
\usepackage{hyperref}
\title{Way\'oro Language Resources for K-12 Curriculum Development}
\author{Compiled by Claude Sonnet 5.5 (claude-sonnet-5-5) for Andr\'as Kornai}
\date{}
\begin{document}
\maketitle

\section{Introduction}
Wayor\'o (ISO 639-3: wyr) is a Tupian language of the Tuparian branch, spoken in Rond\^onia, Brazil. It is critically endangered. Sources I saw report only a handful of elderly speakers, with figures that vary from source to source. The ethnic population is about 250. Documentation began around 2008. Very little exists, and nothing that is a school curriculum was found.

\section{Primary Educational Resources}
No primers, textbooks or school programmes for Wayor\'o were found in these searches.

\section{Linguistic and Reference Materials}
\begin{itemize}
\item General overview: \url{https://en.wikipedia.org/wiki/Wayor%C3%B3_language}.
\item ``Em busca das classes de palavras da l\'{\i}ngua Wayoro'' (word classes in Wayoro), in a journal of the Universidade Federal do Rio de Janeiro: \url{https://revistas.ufrj.br/index.php/rl/article/view/4614/33325}.
\item A paper by Galucio and Nogueira, in the CILLA V conference materials, hosted by AILLA (University of Texas): \url{https://ailla-legacy.lib.utexas.edu/sites/default/files/documents/Galucio_Nogueira_CILLA_V.pdf}. I did not read it, so its exact topic is unconfirmed.
\item Antônia Fernanda de Souza Nogueira and colleagues have published descriptions and a proposed orthography, according to the search summary. I did not confirm specific titles.
\end{itemize}

\section{Religious and Community Resources}
None found.

\section{Digital and Audio Resources}
No audio or video resources were found. The AILLA hosting of the paper above suggests an archive connection, but I did not confirm any Wayor\'o collection.

\section{Curriculum Development Context}
Any curriculum would have to start from scratch with the few remaining speakers and the researchers who work with them. The existing literature is academic and in Portuguese and English.

\end{document}